%! Rational Functions and End Behavior — Unit 3, Lesson 3.1 — Guided Notes
\documentclass[10pt]{article}
\usepackage{precalculus-article}
\usepackage{precalculus-boxes}
\newcommand{\TallMath}[1]{$\displaystyle #1\rule[-1.4em]{0pt}{3.2em}$}

\begin{document}

\pageheader{Unit 3, Lesson 3.1}{Guided Notes}
\namedateperiod

\vspace{0.10in}

\begin{objectivebox}
By the end of this lesson, I will be able to\ldots
\begin{enumerate}[itemsep=2pt]
  \item Determine the end behavior of a rational function by comparing the degrees of the
        numerator and denominator using the \blank{3.5cm}. \hfill\textit{(LO 1.7.A)}
  \item Write the end behavior of a rational function in \blank{2.8cm} notation and identify
        \blank{2.8cm} or \blank{2.4cm} asymptotes. \hfill\textit{(LO 1.7.A)}
\end{enumerate}
\end{objectivebox}

\vspace{0.08in}

\begin{vocabbox}
\termblanklong{rational function}
\termblanklong{quotient of leading terms}
\termblanklong{horizontal asymptote}
\termblanklong{slant asymptote}
\end{vocabbox}

\vspace{0.08in}

\begin{hookbox}
A therapist charges \$200 per session plus a one-time \$50 intake fee.
Average cost per session: $r(n) = \dfrac{200n + 50}{n}$

\medskip
\begin{tabular}{|c|c|c|c|c|}
\hline
Sessions $n$ & 1 & 10 & 100 & 1000 \\
\hline
Avg.\ cost $r(n)$ & \blank{1.8cm} & \blank{1.8cm} & \blank{1.8cm} & \blank{1.8cm} \\
\hline
\end{tabular}

\smallskip
As $n \to \infty$, the average cost approaches \blank{2cm} because \writelines{1}
\end{hookbox}

\vspace{0.08in}

\begin{notesbox}{Quotient of Leading Terms}
\textbf{Key Idea:} For large $|x|$, a polynomial is dominated by its \blank{2.8cm}.
So for a rational function, the end behavior is determined by the \textit{quotient of the leading terms}.

\medskip
\textbf{Worked Example.} Let $r(x) = \dfrac{3x^2 + 1}{x^2 - 4}$.

\medskip
\renewcommand{\arraystretch}{1.5}
\begin{tabular}{@{}p{0.45\linewidth} p{0.45\linewidth}@{}}
Step 1: Leading term of numerator: & \blank{3cm} \\
Step 2: Leading term of denominator: & \blank{3cm} \\
Step 3: Quotient of leading terms: & \blank{3cm} \\
Step 4: Horizontal asymptote: $y = $ & \blank{2cm} \\
\end{tabular}

\medskip
In limit notation:
$\displaystyle\lim_{x\to\infty} r(x) = $ \blank{2.5cm}
\quad and \quad
$\displaystyle\lim_{x\to -\infty} r(x) = $ \blank{2.5cm}

\medskip
\textbf{Numerical check:} $r(100) = \dfrac{30{,}001}{9{,}996} \approx$ \blank{2cm} \quad (close to the HA!)
\end{notesbox}

\vspace{0.08in}

\begin{notesbox}{The Three End-Behavior Cases}
Compare the degree of the numerator (deg num) to the degree of the denominator (deg denom):

\medskip
\renewcommand{\arraystretch}{2.0}
\begin{tabularx}{\linewidth}{@{} >{\bfseries}p{2.4cm} X X @{}}
Case & Degree Relationship & End Behavior \\
\hline
Case 1 & deg num \textbf{$<$} deg denom & \blank{5cm} \\
Case 2 & deg num \textbf{$=$} deg denom & \blank{5cm} \\
Case 3 & deg num \textbf{$>$} deg denom & \blank{5cm} \\
\end{tabularx}

\medskip
\textit{Special case of Case 3:} If deg num $=$ deg denom $+ 1$, the graph has a
\blank{2.5cm} asymptote parallel to a \blank{2cm}.

\medskip
\textbf{Quick Practice.} Classify each rational function below:

\renewcommand{\arraystretch}{1.8}
\begin{tabular}{@{}p{3.8cm} p{1.5cm} p{1.5cm} p{2.5cm} p{2.5cm}@{}}
\textbf{Function} & \textbf{Deg num} & \textbf{Deg den} & \textbf{Case} & \textbf{Asymptote} \\
\hline
$\dfrac{x+1}{x^2-9}$         & \blank{1cm} & \blank{1cm} & \blank{1cm} & \blank{2cm} \\
$\dfrac{3x^2}{x^2+4}$        & \blank{1cm} & \blank{1cm} & \blank{1cm} & \blank{2cm} \\
$\dfrac{x^3-2x}{x-3}$        & \blank{1cm} & \blank{1cm} & \blank{1cm} & \blank{2cm} \\
$\dfrac{2x+5}{x-1}$          & \blank{1cm} & \blank{1cm} & \blank{1cm} & \blank{2cm} \\
\end{tabular}
\end{notesbox}

\end{document}
